% Independent coefficient-three calculation and its exact tower consequence.
\subsection{At coefficient three only the matching obstruction remains}

\begin{result}{E.66}{Proved}{Vanishing of the finite-ramification second group at three in every weight}
\label{coh:threeallweights}
Put $G_S=G_{\mathbb Q,\{3,\infty\}}$. For every integer $w$ and
every $n\geq1$,
\[
 H^2(G_S,\mathbb Z/3^n(w))=0,
 \qquad
 H^2_{\rm cts}(G_S,\mathbb Z_3(w))=0.
\]
In particular, for every $w\geq2$, A.22 specializes to an
isomorphism
\[
 H^2_{\rm cts}(G_{\mathbb Q},\mathbb Z_3(w))
 \ \xrightarrow{\ \sim\ }\
 \frac{c_0(\{q:q\ne3\},\mathbb Z_3)}
      {\{((q^{w-1}-1)y_q)_q:(y_q)_q\in c_0(\mathbb Z_3)\}}.
\]
The denominator is the actual image. This absolute group need
not vanish; the conclusion concerns its finite-ramification
kernel.

Consequently, for the ordinary pairs and boundary matching
specified in B.25, every boundary-matched pro-$3$ class-$r$
section, $r\geq2$, has a global integral class-$(r+1)$ lift.
Such a lift can be chosen boundary matched if and only if the
discrepancy tuple satisfies the $H^1$ localization-image condition
of B.26. Thus at coefficient three that matching condition is
the sole remaining step at every depth.
\end{result}

\begin{proof}
We first compute the finite group with coefficients of order
three. The mod-three cyclotomic character has image
$\mathbb F_3^\times=\{1,-1\}$. Hence
\[
 \mathbb F_3(w)\simeq
 \begin{cases}
   \mathbb F_3(2),&w\text{ even},\\
   \mu_3,&w\text{ odd}.
 \end{cases}
\]
These are isomorphisms of $G_S$-modules, also for zero or negative
$w$.

For the even case, the integral calculation used in A.25 gives
\[
 H^2_{\rm cts}(G_S,\mathbb Z_3(2))
   =K_2(\mathbb Z[1/3])\{3\}=0.
\]
Indeed $K_2(\mathbb Z)=\mathbb Z/2$, and inverting three adds
only a possible residue in
$K_1(\mathbb F_3)$, whose order is two
\cite[III.5.2.2; IV.1.13; V.6.8; VI.8.2]{CohWeibelKbook}.
Also $H^3_{\rm cts}(G_S,\mathbb Z_3(2))=0$.
To retain the meaning of this last assertion, finite three-primary
cohomology of $G_S$ vanishes above degree two and is finite in
degree two \cite[I.4.10(a),(c), pp.56--57]{CohMilneADT}.
Its degree-two inverse system is therefore Mittag--Leffler;
the continuous-cohomology inverse-limit exact sequence gives
the asserted integral degree-three vanishing.
The coefficient sequence
\[
 0\longrightarrow\mathbb Z_3(2)
   \xrightarrow{\ 3\ }\mathbb Z_3(2)
   \longrightarrow\mathbb F_3(2)\longrightarrow0
\]
now gives $H^2(G_S,\mathbb F_3(2))=0$.
In particular, this step uses integral degree two and degree
three, rather than rational vanishing.

Here is the odd case directly in finite coefficients. Purity
and localization on $\operatorname{Spec}\mathbb Z[1/3]$ give
the exact segment
\[
 H^1(G_{\mathbb Q},\mu_3)
 \longrightarrow \bigoplus_{q\ne3}\mathbb F_3
 \longrightarrow H^2(G_S,\mu_3)
 \xrightarrow{\rm inf} H^2(G_{\mathbb Q},\mu_3).
\]
The precise comparison with $G_S$ cohomology is
\cite[II.2.9, pp.170--171]{CohMilneADT}, and the supports
calculation is \cite[II.1.7(b), pp.151--152]{CohMilneADT}.
The first map, under
$H^1(G_{\mathbb Q},\mu_3)=\mathbb Q^\times/\mathbb Q^{\times3}$,
is $(v_q\bmod3)_{q\ne3}$. It is surjective: a finitely
supported vector is represented by the finite product of
the corresponding rational primes to the chosen exponents.
Thus the inflation map in the displayed sequence is injective.
Kummer theory identifies its target with
$\operatorname{Br}(\mathbb Q)[3]$.

An inflated class is zero at every $q\ne3$, since it factors
locally through the unramified quotient and a finite field has
zero second cohomology with these torsion coefficients.
Its real restriction is zero, since a group of order two has
zero positive-degree cohomology on a three-primary module.
The global Brauer injection and the sum of local invariants
then force its three-adic restriction to be zero as well:
only that one local invariant could remain, and their sum is
zero. The global class is zero. The exact sequence, including
this form of Brauer reciprocity, is stated and used in
\cite[II.2.1--2.2, pp.163--165]{CohMilneADT}.
Injectivity of inflation proves
$H^2(G_S,\mu_3)=0$. We have proved
$H^2(G_S,\mathbb F_3(w))=0$ for every integer $w$.

For clarity, passing to higher finite exponents uses the
following literal short exact sequence. Write
$T=\mathbb Z_3(w)$ and $T_n=T/3^nT$. Define
\[
 0\longrightarrow T_1
 \xrightarrow{\ i_n\ }T_{n+1}
 \xrightarrow{\ \pi_n\ }T_n\longrightarrow0,
 \qquad
 i_n(t\bmod3T)=3^nt\bmod3^{n+1}T .
\]
The map $i_n$ is well-defined and injective, its image is the
kernel of reduction $\pi_n$, and both maps are $G_S$-equivariant.
Equivariance follows from multiplication by
$\chi_3(g)^w$ on the fixed Tate module $T$; no identification
between twists at different exponents is being assumed.
The cohomology sequence contains
\[
 H^2(G_S,T_1)\longrightarrow H^2(G_S,T_{n+1})
                  \longrightarrow H^2(G_S,T_n).
\]
The outer groups vanish by the order-three case and induction.
Therefore every middle group vanishes. This induction requires
no assertion about a Bockstein or a third cohomology group at
the higher exponents.

It remains to pass to continuous integral coefficients. Each
$H^1(G_S,T_n)$ is finite. Indeed a cocycle defines a continuous
homomorphism into the finite semidirect product of $T_n$ by
the finite image of $\chi_3^w\bmod3^n$. The resulting finite
Galois extension has bounded degree and ramifies only at three
and infinity. The fixed-ramification finiteness theorem
\cite[Theorem 8.42, p.151]{CohMilneANT} gives finitely many
such extensions, and each allows finitely many homomorphisms.
Thus the inverse system of first cohomology groups is
Mittag--Leffler. Jannsen's exact sequence
\cite[(2.1) and Theorem 2.2, p.215]{CohJannsen1988} is
\[
 0\longrightarrow\varprojlim\nolimits^1H^1(G_S,T_n)
 \longrightarrow H^2_{\rm cts}(G_S,T)
 \longrightarrow\varprojlim_n H^2(G_S,T_n)\longrightarrow0.
\]
Both outside terms are zero, by finiteness on the left and
the finite calculation on the right. This proves the
continuous assertion separately from the finite one.

The strict localization description is now exactly A.22 with
its first term zero. For the final tower assertion, B.25
puts the next obstruction of a boundary-matched class-$r$
section in
$H^2_{\rm cts}(G_S,\mathbb Z_3(r+1))^{b_{r+1}}$.
That group is zero, so a continuous integral global lift
exists. B.26 describes precisely which of its global
$H^1$ twists match the local point classes. It remains an
independent localization-image test, with a finite-place,
finite-coefficient witness whenever it fails.
\end{proof}

In particular, odd weights and integral three-power torsion
are included in E.66. No inference from rational
cohomology or from a nonvanishing special value is used.
The conclusion removes the finite arithmetic gate in B.25
at three; it does not supply a compatible sequence of
boundary-matched lifts for any infinite-support pair.

\begin{result}{E.67}{Proved}{The exact integral first-cohomology parameters at three}
\label{coh:threeh1}
For every $w\geq2$, inflation identifies
$H^1(G_{\mathbb Q},\mathbb Z_3(w))$ with
$H^1(G_S,\mathbb Z_3(w))$, and as compact
$\mathbb Z_3$-modules
\[
 H^1(G_{\mathbb Q},\mathbb Z_3(w))\simeq
 \begin{cases}
   \mathbb Z/3^{\,1+v_3(w)},&w\text{ even},\\
   \mathbb Z_3,&w\text{ odd}.
 \end{cases}
\]
For even $w$, put $s_w=1+v_3(w)$. A generator is the
integral cyclotomic cocycle
\[
 \eta_w(g)=\frac{\chi_3(g)^w-1}{3^{s_w}}.
\]
At each $q\ne3$, its localization is the class of
$(q^w-1)/3^{s_w}$ in
$\mathbb Z_3/(q^w-1)\mathbb Z_3$ and has exact order
$3^{s_w}$. In particular, localization of this global
group at any one prime $q\ne3$ is injective.

At a weight-$w$ step in B.26, the global matching parameters
are therefore exactly
$(\mathbb Z/3^{s_w})^{b_w}$ when $w$ is even, and a compact
free module $\mathbb Z_3^{b_w}$ when $w$ is odd. At an even
step, if there is a support prime different from three,
there is at most one boundary-matched global lift class
above a given preceding class.
\end{result}

\begin{proof}
The tame relation in A.20 shows that every integral
weight-$w$ cocycle is unramified away from three.
The same statement for coboundaries is immediate.
This proves the inflation assertion. It also shows,
by A.20, that the first-cohomology group $H$ is a compact
finitely generated $\mathbb Z_3$-module: the cocycle
module is finite free, and the coboundary submodule is
the compact image of $\mathbb Z_3(w)$, hence is closed.

We compute the finite first group with coefficients
$\mathbb F_3(w)$. If $w$ is odd, it is
$H^1(G_S,\mu_3)$. Inflation into
$\mathbb Q^\times/\mathbb Q^{\times3}$ identifies it
with the kernel of the valuations modulo three away
from three, by the first part of the same finite
purity sequence used in E.66. A rational number with
all those valuations divisible by three is a power
of three times a rational cube; its sign is also a
cube. Thus this group is one-dimensional over
$\mathbb F_3$, generated by the Kummer class of three.

If $w$ is even, the module is the trivial module
$\mathbb F_3$, so the group consists of the continuous
homomorphisms $G_S\to\mathbb F_3$. Each nonzero one
defines a cyclic cubic extension of $\mathbb Q$
unramified outside three. Kronecker--Weber
\cite[I.4.16 and its proof, pp.49--50]{RootMilneCFT}
puts such an extension in a cyclotomic field.
The inertia factors at primes other than three must
act trivially on the cubic quotient. Equivalently,
the proof of that theorem gives a containing
cyclotomic field of three-power conductor when
three is the only ramified prime.
The maximal pro-three quotient of
$\operatorname{Gal}(\mathbb Q(\mu_{3^\infty})/\mathbb Q)
=\mathbb Z_3^\times$ is
$1+3\mathbb Z_3\simeq\mathbb Z_3$.
It has exactly a one-dimensional space of
homomorphisms to $\mathbb F_3$; a nonzero one is
already visible in the cyclic cubic subfield of
$\mathbb Q(\mu_9)$. This proves
\[
 \dim_{\mathbb F_3}H^1(G_S,\mathbb F_3(w))=1
 \quad\text{in both parities}.
\]

Apply the coefficient exact sequence for multiplication
by three on $T=\mathbb Z_3(w)$. E.66 gives
$H^2_{\rm cts}(G_S,T)[3]=0$, so it gives an isomorphism
\[
 H/3H \ \xrightarrow{\ \sim\ }\
 H^1(G_S,\mathbb F_3(w)).
\]
Hence the compact finitely generated module $H$ is
cyclic by Nakayama's lemma. Explicitly, if $h$ lifts
a generator modulo three, successively correcting
by its multiples modulo $3^n$ expresses every
element as a convergent $\mathbb Z_3$ multiple of $h$.
Thus $H$ is either $\mathbb Z_3$ or a nonzero finite
cyclic three-group.

We determine its torsion directly, without a rank
theorem for rational cohomology. There are no
invariants in $\mathbb Q_3(w)$, since the cyclotomic
character has infinite image and $w\ne0$. The map
\[
 (\mathbb Q_3(w)/\mathbb Z_3(w))^{G_S}
   \longrightarrow H,\qquad
 v+\mathbb Z_3(w)\longmapsto[g\mapsto g(v)-v]
\]
is therefore injective. It is well-defined because
invariance of the coset makes the cocycle integral;
changing $v$ by an integral element changes it by a
coboundary. Its image is exactly the torsion of $H$.
Indeed its values are killed by a power of three
after clearing the denominator of $v$. Conversely,
if $3^a[c]=0$, choose $t\in T$ with
$3^ac=\delta t$; then
$c=\delta(t/3^a)$, and $t/3^a+T$ is invariant.

For odd $w$, complex conjugation acts as $-1$,
so an invariant coset in this three-primary group
is killed by two and must be zero. Thus $H$ is
torsion free and cyclic with nonzero reduction,
and is $\mathbb Z_3$.

For even $w$, the ideal generated by all
$u^w-1$ with $u\in\mathbb Z_3^\times$ is
$3^{s_w}\mathbb Z_3$. To see the valuation exactly,
write $u$ as its sign times an element of
$1+3\mathbb Z_3$. The sign has trivial $w$th
power. The element $4$ topologically generates
$1+3\mathbb Z_3$, and
$v_3(4^w-1)=1+v_3(w)$ by the binomial expansion
(or the three-adic logarithm). All other elements
give valuations at least this large.
The invariant cosets are consequently precisely
$3^{-s_w}\mathbb Z_3/\mathbb Z_3$, a cyclic group
of order $3^{s_w}$. Since a free cyclic
$\mathbb Z_3$-module has no nonzero torsion, the
already cyclic module $H$ must equal this torsion
subgroup. Its indicated generator maps to
$\eta_w$, proving the formula.

Finally, integral local weight-$w$ cocycles at
$q\ne3$ are unramified by the tame relation.
Evaluation at arithmetic Frobenius identifies their
cohomology with
$\mathbb Z_3/(q^w-1)\mathbb Z_3$.
The displayed value of $\eta_w$ follows from
$\chi_3(\operatorname{Frob}_q)=q$.
Since $v_3(q^w-1)\geq s_w$, this class has
exact order $3^{s_w}$. The final statements
follow by taking $b_w$ copies and applying
the torsor description in B.26.
\end{proof}

\begin{result}{E.68}{Proved}{Matching at the coefficient prime three}
\label{coh:threelocalization}
For $w\geq2$, put
$H_{3,w}=H^1(G_{\mathbb Q_3},\mathbb Z_3(w))$.
This is a compact module of rank one. Its torsion is zero
for odd $w$, and is a cyclic group of order
$3^{s_w}$, $s_w=1+v_3(w)$, for even $w$.
The localization map from the global group in E.67 is an
isomorphism onto $H_{3,w}$ for odd $w$, and an isomorphism
onto its torsion subgroup for even $w$.

Suppose that a global class-$(w-1)$ section agrees, up to
geometric conjugacy, with a prescribed local point section
at three, and that a global class-$w$ lift exists. For odd
$w\geq3$, precisely one global lift class above that
preceding class matches this local point at class $w$.
For even $w$, such a matching lift exists precisely when
the local discrepancy maps to zero in
$H_{3,w}^{b_w}/(H_{3,w}^{\rm tors})^{b_w}$; if it exists,
it is unique. These assertions concern matching at three.
The conditions at the other support primes remain to be
checked for that uniquely specified lift.
\end{result}

\begin{proof}
We give the integral local calculation, including the
finite groups that control the rank. For any nonzero
integer $j$, define
\[
 a_j=
 \begin{cases}
  0,&j\text{ odd},\\
  1+v_3(j),&j\text{ even}.
 \end{cases}
\]
The local cyclotomic character at three has image
$\mathbb Z_3^\times$. The same computation with $-1$
and $4$ as in E.67 gives
\[
 \left|H^0(G_{\mathbb Q_3},\mathbb Z/3^n(j))\right|
       =3^{\min(n,a_j)}.
\]
The compatible finite local duality pairings identify
$H^2(G_{\mathbb Q_3},\mathbb Z/3^n(w))$
with the dual of
$H^0(G_{\mathbb Q_3},\mathbb Z/3^n(1-w))$
\cite[I.2.3, p.28]{CohMilneADT}.
The local Euler characteristic is $3^{-n}$
\cite[I.2.8, p.31]{CohMilneADT}. Consequently
\[
 \left|H^1(G_{\mathbb Q_3},\mathbb Z/3^n(w))\right|
       =3^{\,n+\min(n,a_w)+\min(n,a_{1-w})}.
\]

All these finite local cohomology groups are finite.
Jannsen's inverse-limit sequence therefore has no
derived-limit term here. Dualizing the degree-two
inverse limit identifies
\[
 H^2_{\rm cts}(G_{\mathbb Q_3},\mathbb Z_3(w))
 \simeq
 H^0(G_{\mathbb Q_3},
          (\mathbb Q_3/\mathbb Z_3)(1-w))^\vee .
\]
The invariant group on the right has order $3^{a_{1-w}}$.
The transition maps are essential: the duals of
coefficient reductions are the inclusions of
$3^n$-torsion in $(\mathbb Q_3/\mathbb Z_3)(1-w)$.
Thus this is the actual integral local group.
The coefficient exact sequence now gives
\[
 \left|H_{3,w}/3^nH_{3,w}\right|
       =3^{\,n+\min(n,a_w)}.
\]
Indeed its cokernel in finite first cohomology is
exactly the $3^n$-torsion of the preceding integral
second group.

The module $H_{3,w}$ is compact, as the inverse limit
of its finite first cohomology groups; finiteness of
the degree-zero groups removes the degree-one
derived-limit term. Its quotient modulo three is
finite. The compact Nakayama argument by successive
correction makes it finitely generated.
The displayed growth formula shows that its free
rank is one. Its torsion is computed directly by
the invariant-coset argument of E.67, now over
$G_{\mathbb Q_3}$:
\[
 H_{3,w}^{\rm tors}
   =(\mathbb Q_3(w)/\mathbb Z_3(w))^{G_{\mathbb Q_3}}.
\]
This is zero for odd $w$ and cyclic of order
$3^{s_w}$ for even $w$.

For odd $w$, both the global and local integral first
groups are free of rank one. Reduction of the global
group modulo three is generated by $\kappa(3)$
under the identification
$\mathbb F_3(w)\simeq\mu_3$ used in E.67.
Its local restriction is nonzero: the valuation of
three is one, so three is not a cube in
$\mathbb Q_3^\times$. The natural map
$H_{3,w}/3H_{3,w}\to
H^1(G_{\mathbb Q_3},\mathbb F_3(w))$ is injective
by the coefficient sequence.
Thus localization on the two free rank-one modules
is nonzero modulo three. It is multiplication by
a three-adic unit after choices of bases, hence an
isomorphism.

For even $w$, the explicit generator $\eta_w$
is obtained globally and locally from the same
invariant coset $3^{-s_w}+\mathbb Z_3(w)$.
It therefore generates the local torsion as well.
This proves the assertion about localization.

Finally, the stabilizer calculation in B.26 holds
at three too: on a first nonzero weight-$i$ coordinate
of a putative geometric stabilizer it gives
$(\chi_3(g)^i-1)k_i=0$, and the local cyclotomic
character has infinite image. Thus $k_i=0$, a
contradiction. Alignment with the prescribed local
point below the central extension is unique.
Its lift discrepancy is consequently a well-defined
element of $H_{3,w}^{b_w}$, and changing a global
lift changes it by the global localization map.
The descriptions just proved give existence and
uniqueness in the two parities, exactly as stated.
\end{proof}
